% ---------------------------------------------------------------------
%  Chapitre 5 — Le protocole, ses limites, ses suites (aperçu, rédigé)
% ---------------------------------------------------------------------
\chapter{Le protocole, ses limites, ses suites}
\label{chap:protocole}

\section{Le protocole récapitulé}
\label{sec:protocole}

Ce que le stage livre n'est pas une stratégie gagnante, mais une
procédure de jugement, réutilisable telle quelle sur tout signal futur
de l'agent, qu'il soit technique, informationnel ou relationnel. Le
tableau~\ref{tab:protocole} récapitule cette procédure en six étages.
Chaque étage répare une menace précise identifiée au fil du mémoire, de
la dérive du marché à l'autocorrélation résiduelle, et le verdict n'est
prononcé que si tous les étages concordent. C'est cette exigence
conjointe qui fait la valeur de l'ensemble~: la défaillance d'une seule
condition suffit, comme on l'a vu, à fabriquer un faux positif.

\begin{table}[H]
\centering
\caption{Le protocole de validation en six étages~: chaque étage répare
une menace identifiée, et le verdict exige de tous les franchir.}
\label{tab:protocole}
\small
\begin{tabular}{@{}p{0.30\linewidth}p{0.34\linewidth}p{0.26\linewidth}@{}}
\toprule
Étage & Menace traitée & Outil \\
\midrule
1. \edgevar{} apparié au hasard & dérive du marché & Monte Carlo apparié \\
2. Student unilatéral + TCL & variabilité d'échantillonnage & test de référence \\
3. Seuil de Bonferroni & multiplicité (11 tests) & contrôle du risque global \\
4. DEFF + bootstrap par grappes & dépendance intra-titre & sondages, rééchantillonnage \\
5. Les deux portes (titre, mois) & pseudo-réplication & agrégation équipondérée \\
6. Série mensuelle AR($p$) & autocorrélation résiduelle & variance de long terme \\
\bottomrule
\end{tabular}
\end{table}

\section{Limites assumées}
\label{sec:limites}

Aucune des limites qui suivent n'invalide le verdict rendu. Chacune
trace en revanche la frontière de ce que le protocole peut affirmer.

\begin{limites}
\begin{itemize}
  \item Résultats établis sur l'échantillon ayant servi aux calculs~:
        une validation hors échantillon (\emph{walk-forward}) reste
        nécessaire avant tout usage opérationnel.
  \item Périmètre du jugement~: dix stratégies déclarées à l'avance. À
        plus grande échelle, le contrôle de la fouille de données
        exigerait des outils dédiés comme le \emph{Reality Check} de
        White~\cite{white2000}.
  \item Frais et conditions d'exécution modélisés de façon simplifiée.
  \item Puissance élevée des grands échantillons~: distinguer
        soigneusement significativité et taille d'effet.
  \item Biais du survivant~: la base ne contient que les entreprises
        présentes dans l'indice sur la période, les sociétés disparues
        en étant absentes~; ce biais est atténué par la construction
        appariée de l'\edgevar{}, qui compare chaque trade au hasard sur
        le même titre.
\end{itemize}
\end{limites}

\section{Perspectives}
\label{sec:perspectives}

La première suite naturelle est la validation temporelle. Elle consiste
à figer les règles sur une période d'apprentissage puis à juger sur la
période suivante, en avançant pas à pas, ce que la littérature appelle
le walk-forward. La validation croisée classique est ici interdite,
puisqu'elle mélangerait le futur et le passé.

La deuxième est l'avantage conditionnel. Les cartes du
chapitre~\ref{chap:contexte} suggèrent que le succès dépend du régime de
marché~; modéliser la probabilité de gain en fonction du contexte, par
exemple par une régression logistique, en serait le prolongement direct.
Cette voie a été écartée du présent mémoire par rigueur~: sur des
dizaines de milliers de trades dépendants, les p-values des coefficients
seraient trop optimistes, exactement le piège du
chapitre~\ref{chap:juger}. Les outils adaptés, modèles à effets mixtes
ou équations d'estimation généralisées, dépassent le cadre de ce travail
et en constituent la prochaine étape. Une version outillée du protocole
devrait aussi intégrer un garde-fou d'effectif, écartant en amont les
stratégies trop peu fournies pour que le théorème central limite
s'applique.

La troisième concerne les autres signaux de l'agent. Celui-ci collecte
également des signaux d'information, contrats publics, transactions de
responsables politiques, annonces réglementaires, ainsi que des
relations d'avance ou de retard entre titres, au sens de la causalité de
Granger~\cite{granger1969}. Des travaux exploratoires ont été engagés
sur ces deux fronts. Le protocole du chapitre~\ref{chap:juger} est
précisément l'outil qui permettra de les juger, hors du périmètre de ce
mémoire.
